\begin{titlepage}
\thispagestyle{empty}
    \begin{center}
        \begin{spacing}{1.15}
        \textbf{\large RANCANG BANGUN INSTRUMENTASI SENSOR \textit{TUNNELING MAGNETORESISTANCE} BERBASIS NANOFIBER \ch{Fe3O4}/PVA-SITRAT-ADH UNTUK DETEKSI FORMALIN PADA BAKSO MENGGUNAKAN KOMPARASI MODEL KLASIK (SVM, RANDOM FOREST) DAN KUANTUM (QSVC, VQC)}
        \end{spacing}
                
        \vfill
        \textbf{\large PROPOSAL TUGAS AKHIR}
        
        \vspace{0.4cm}
        \begin{spacing}{1.1}
        \textbf{Diajukan Sebagai Salah Satu Syarat Melaksanakan Penelitian Tugas Akhir \\ Jurusan Fisika Fakultas Sains dan Teknologi \\ Universitas Islam Negeri Sunan Gunung Djati Bandung}
        \end{spacing}

        \vfill
        \includegraphics[width=4.2cm]{Logo UIN.png}

        \vfill
        \begin{spacing}{1.1}
        \textbf{\underline{Fahry Rizky Samsudin}\\
        1237030018}
        \end{spacing}
        
        \vfill
        \begin{spacing}{1.15}
        \textbf{JURUSAN FISIKA\\
        FAKULTAS SAINS DAN TEKNOLOGI\\
        UNIVERSITAS ISLAM NEGERI SUNAN GUNUNG DJATI BANDUNG\\
        \vspace{0.3cm}
        \large 2026}
        \end{spacing}
    \end{center}
\end{titlepage}